% !TEX root = ../LosningsforslagTRES0312.tex
\chapter{Termofysikk}
\label{LF:Chapter:termofysikk}
Dette kapittelet inneholder løsningsforslag til oppgavene i Kapittel~\ref{komp-Chapter:termofysikk}, \textit{Termofysikk}, i kompendiet.

\oppgaver{Omregning av temperaturer}{k8oppg9}{%
%
Gjør om temperaturene i a)--d).
\begin{enumerate}[a)]
    \item Normal kroppstemperatur, $37\,^\circ$C, til kelvin.
    \item Kokepunktet til flytende nitrogen, \SI{77}{\kelvin}, til celsius.
    \item Temperaturen på overflaten av sola, \SI{5.8e3}{\kelvin}, til celsius.
    \item En amerikansk værmelding lover $50\,^\circ$F. Hva er det i celsius?
    \item En morgen er det $-5\,^\circ$C, og om ettermiddagen er det $15\,^\circ$C. Hvor stor er temperaturendringen $\Delta T$ i kelvin?
    \item Ved hvilken temperatur viser en celsius-skala og en fahrenheit-skala det samme tallet?
\end{enumerate}
%
%
}

\svar[title={Løsningsforslag}]{%
\noindent
Vi bruker $T = t + 273$ og $t_\text{F} = \tfrac{9}{5}\,t_\text{C} + 32$.

\smallskip\noindent
\textbf{a)} $T = 37 + 273 = \doubleunderline{\SI{310}{\kelvin}}$

\smallskip\noindent
\textbf{b)} $t = 77 - 273 = \doubleunderline{-196\,^\circ\text{C}}$

\smallskip\noindent
\textbf{c)} $t = 5800 - 273 = 5527\,^\circ\text{C} \approx \doubleunderline{5.5\cdot 10^3\,^\circ\text{C}}$. Når temperaturen er så høy, spiller det liten rolle om vi bruker kelvin eller celsius.

\smallskip\noindent
\textbf{d)} Vi løser formelen for $t_\text{C}$:
\[
  t_\text{C} = \frac{5}{9}\,(t_\text{F} - 32) = \frac{5}{9}\,(50 - 32) = \doubleunderline{10\,^\circ\text{C}}
\]

\smallskip\noindent
\textbf{e)} En grad celsius er like stor som én kelvin, så
\[
  \Delta T = 15 - (-5) = \doubleunderline{\SI{20}{\kelvin}}
\]

\smallskip\noindent
\textbf{f)} Vi setter $t_\text{F} = t_\text{C} = t$:
\[
  t = \frac{9}{5}\,t + 32 \quad\implies\quad -\frac{4}{5}\,t = 32 \quad\implies\quad t = \doubleunderline{-40}
\]
Altså er $-40\,^\circ$C det samme som $-40\,^\circ$F.
}

\oppgaver{Dobbelt så varmt?}{k8oppg10}{%
%
\begin{enumerate}[a)]
    \item En dag stiger temperaturen fra $10\,^\circ$C til $20\,^\circ$C. Har den gjennomsnittlige kinetiske energien til luftmolekylene blitt dobbelt så stor? \textit{(Hint: Den gjennomsnittlige kinetiske energien er proporsjonal med temperaturen målt i kelvin.)}
    \item En morgen er det $-5\,^\circ$C, og om ettermiddagen er det $15\,^\circ$C. Hvor stor er temperaturendringen $\Delta T$ i kelvin?
    \item Ved hvilken temperatur viser en celsius-skala og en fahrenheit-skala det samme tallet?
\end{enumerate}
%
%
}

\svar[title={Løsningsforslag}]{%
\noindent
Den gjennomsnittlige kinetiske energien er proporsjonal med den absolutte temperaturen, $\bar{E}_\text{K} = \tfrac{3}{2}k_\text{B}T$. Vi må derfor regne i kelvin: $T_1 = 10 + 273 = \SI{283}{\kelvin}$ og $T_2 = 20 + 273 = \SI{293}{\kelvin}$. Da er
\[
  \frac{\bar{E}_{\text{K},2}}{\bar{E}_{\text{K},1}} = \frac{\tfrac{3}{2}k_\text{B}T_2}{\tfrac{3}{2}k_\text{B}T_1} = \frac{T_2}{T_1} = \frac{\SI{293}{\kelvin}}{\SI{283}{\kelvin}} = 1.035 \approx \doubleunderline{1.04}
\]
Energien er altså bare omtrent \SI{3.5}{\percent} større, selv om temperaturen i celsius er dobbelt så høy. Det er bare i kelvin det gir mening å si «dobbelt så høy temperatur».
}

\oppgaver{Oksygenmolekyler}{k8oppg11}{%
%
\begin{enumerate}[a)]
    \item Finn den gjennomsnittlige kinetiske energien til et oksygenmolekyl ved $0\,^\circ$C.
    \item Finn den gjennomsnittlige kinetiske energien ved $100\,^\circ$C.
    \item Hvor mange prosent større er energien ved $100\,^\circ$C enn ved $0\,^\circ$C?
\end{enumerate}
%
%
}

\svar[title={Løsningsforslag}]{%
\noindent
Vi bruker $\bar{E}_\text{K} = \tfrac{3}{2}k_\text{B}T$ med temperaturen i kelvin.

\smallskip\noindent
\textbf{a)} $T = 0 + 273 = \SI{273}{\kelvin}$:
\[
  \bar{E}_\text{K} = \frac{3}{2}\cdot\SI{1.38e-23}{\joule\per\kelvin}\cdot\SI{273}{\kelvin} = \doubleunderline{\SI{5.65e-21}{\joule}}
\]

\smallskip\noindent
\textbf{b)} $T = 100 + 273 = \SI{373}{\kelvin}$:
\[
  \bar{E}_\text{K} = \frac{3}{2}\cdot\SI{1.38e-23}{\joule\per\kelvin}\cdot\SI{373}{\kelvin} = \doubleunderline{\SI{7.72e-21}{\joule}}
\]

\smallskip\noindent
\textbf{c)}
\[
  \frac{\SI{7.72e-21}{\joule}}{\SI{5.65e-21}{\joule}} = \frac{373}{273} = 1.37
\]
Energien er \doubleunderline{\SI{37}{\percent}} større. Legg merke til at svaret ikke avhenger av hva slags gass det er.
}

\oppgaver{Helium og oksygen}{k8oppg12}{%
%
En beholder inneholder en blanding av heliumgass og oksygengass med temperatur \SI{300}{\kelvin}. Massen til et heliumatom er \SI{6.64e-27}{\kilo\gram}, og massen til et oksygenmolekyl er \SI{5.31e-26}{\kilo\gram}.
\begin{enumerate}[a)]
    \item Hvilken av de to har størst gjennomsnittlig kinetisk energi?
    \item Bruk $\bar{E}_\text{K} = \tfrac{1}{2}mv^2$ til å finne en typisk fart for heliumatomene og for oksygenmolekylene.
    \item Hvor mange ganger raskere beveger heliumatomene seg?
\end{enumerate}
%
%
}

\svar[title={Løsningsforslag}]{%
\noindent
\textbf{a)} Den gjennomsnittlige kinetiske energien avhenger bare av temperaturen, $\bar{E}_\text{K} = \tfrac{3}{2}k_\text{B}T$. Begge gassene har samme temperatur, og har derfor \doubleunderline{like stor} gjennomsnittlig kinetisk energi:
\[
  \bar{E}_\text{K} = \frac{3}{2}\cdot\SI{1.38e-23}{\joule\per\kelvin}\cdot\SI{300}{\kelvin} = \SI{6.21e-21}{\joule}
\]

\smallskip\noindent
\textbf{b)} Vi løser $\bar{E}_\text{K} = \tfrac{1}{2}mv^2$ for $v$:
\[
  v = \sqrt{\frac{2\bar{E}_\text{K}}{m}}
\]
For helium:
\[
  v_\text{He} = \sqrt{\frac{2\cdot\SI{6.21e-21}{\joule}}{\SI{6.64e-27}{\kilo\gram}}} = \doubleunderline{\SI{1.37e3}{\metre\per\second}}
\]
For oksygen:
\[
  v_{\text{O}_2} = \sqrt{\frac{2\cdot\SI{6.21e-21}{\joule}}{\SI{5.31e-26}{\kilo\gram}}} = \doubleunderline{\SI{484}{\metre\per\second}}
\]

\smallskip\noindent
\textbf{c)}
\[
  \frac{v_\text{He}}{v_{\text{O}_2}} = \frac{\SI{1.37e3}{\metre\per\second}}{\SI{484}{\metre\per\second}} = \doubleunderline{2.83}
\]
Heliumatomene er lettere, og må derfor bevege seg raskere for å ha like stor kinetisk energi.
}

\oppgaver{Riktig eller galt?}{k8oppg2}{%
%
Er påstandene riktige eller gale? Begrunn svaret.
\begin{enumerate}[a)]
    \item En kopp varm kaffe inneholder mye varme.
    \item Varme går alltid fra legemet med mest indre energi til legemet med minst indre energi.
    \item To legemer som har samme temperatur, utveksler ikke varme med hverandre.
    \item Temperaturen til et stoff kan øke uten at vi tilfører varme.
\end{enumerate}
%
%
}

\svar[title={Løsningsforslag}]{%
\noindent
\textbf{a)} \doubleunderline{Galt}. Varme er energi \emph{på vei} mellom to systemer. Kaffen har mye \emph{indre energi}, men den «inneholder» ikke varme. Den kan derimot \emph{avgi} varme til en kaldere omgivelse.

\smallskip\noindent
\textbf{b)} \doubleunderline{Galt}. Varme går alltid fra høy \emph{temperatur} til lav temperatur. Et badekar med lunkent vann har mye mer indre energi enn en kopp kokende vann, men varmen går likevel fra koppen til badevannet hvis vi setter koppen i karet.

\smallskip\noindent
\textbf{c)} \doubleunderline{Riktig}. Varme strømmer bare når det er en temperaturforskjell. To legemer med samme temperatur er i \emph{termisk likevekt}.

\smallskip\noindent
\textbf{d)} \doubleunderline{Riktig}. Vi kan også øke den indre energien ved å gjøre \emph{arbeid} på stoffet, for eksempel med friksjon eller ved å presse sammen en gass, som i en sykkelpumpe.
}

\oppgaver{Les av oppvarmingskurven}{k8oppg5}{%
%
Bruk oppvarmingskurven for \SI{1}{\kilo\gram} vann i figur~\komplabel{fig:fasediagram_vann}.
\begin{enumerate}[a)]
    \item Hva skjer med temperaturen mens isen smelter? Hvor mye energi går med til å smelte isen?
    \item Omtrent hvor mye energi trengs for å varme opp vannet fra $0\,^\circ$C til $100\,^\circ$C?
    \item Hvorfor er linjen for is brattere enn linjen for vann?
    \item Omtrent hvor mye energi trengs totalt for å gjøre isen ved $-20\,^\circ$C om til damp ved $100\,^\circ$C?
\end{enumerate}
%
%
}

\svar[title={Løsningsforslag}]{%
\noindent
\textbf{a)} Mens isen smelter, er temperaturen \doubleunderline{konstant, $0\,^\circ$C}. Det er det horisontale platået i figuren. Energien som går med, er \doubleunderline{\SI{334}{\kilo\joule}}. Energien brukes til å bryte bindingene i iskrystallen, ikke til å gi molekylene større fart.

\smallskip\noindent
\textbf{b)} Linjen for vann starter ved omtrent \SI{376}{\kilo\joule} og slutter ved omtrent \SI{796}{\kilo\joule}. Energien er differansen:
\[
  E \approx \SI{796}{\kilo\joule} - \SI{376}{\kilo\joule} = \doubleunderline{\SI{4.2e2}{\kilo\joule}}
\]

\smallskip\noindent
\textbf{c)} Is har mindre spesifikk varmekapasitet enn vann (omtrent halvparten). Det trengs derfor mindre energi for å øke temperaturen til isen med én grad, og temperaturen stiger raskere per kilojoule. Linjen blir brattere.

\smallskip\noindent
\textbf{d)} Fordampingen er ferdig der det øverste platået slutter, ved omtrent \SI{3055}{\kilo\joule}. Totalt trengs det altså omtrent \doubleunderline{\SI{3.1}{\mega\joule}}. Legg merke til at fordampingen alene står for \SI{2259}{\kilo\joule}, altså mer enn to tredjedeler av energien.
}

\oppgaver{Hvorfor svetter vi?}{k8oppg6}{%
%
\begin{enumerate}[a)]
    \item Når svette fordamper fra huden, trengs det energi. Hvor kommer denne energien fra, og hva skjer med kroppen?
    \item Hvorfor føles det kaldt å komme opp av vannet på en dag med vind, selv om lufta er varm?
    \item Hvorfor gir vanndamp på $100\,^\circ$C verre brannskader enn vann på $100\,^\circ$C?
\end{enumerate}
%
%
}

\svar[title={Løsningsforslag}]{%
\noindent
\textbf{a)} Fordamping krever energi (fordampingsvarme). Energien tas fra huden og kroppen, som avgir varme til svetten. Kroppen blir dermed \doubleunderline{avkjølt}. Det er derfor vi svetter når vi er varme.

\smallskip\noindent
\textbf{b)} Når du er våt, fordamper vannet fra huden, og fordampingen tar energi fra kroppen. Vinden fjerner den fuktige lufta like ved huden, slik at vannet fordamper \emph{raskere}. Da tas det mer energi fra kroppen per sekund, og du fryser, selv om lufta er varm.

\smallskip\noindent
\textbf{c)} Vi sammenligner per kilogram og leser av figuren.

\emph{Vann på $100\,^\circ$C} avgir bare varme mens det kjøles ned til $37\,^\circ$C. Fra figuren trengs det omtrent \SI{420}{\kilo\joule} for å varme \SI{1}{\kilo\gram} vann \SI{100}{\kelvin}. Avkjøling på \SI{63}{\kelvin} gir derfor omtrent
\[
  \frac{63}{100}\cdot\SI{420}{\kilo\joule} \approx \SI{2.6e2}{\kilo\joule}
\]
\emph{Damp på $100\,^\circ$C} kondenserer først til vann og avgir da like mye energi som det trengs for å fordampe det, \SI{2259}{\kilo\joule}. Deretter kjøles vannet ned som over. Til sammen er det omtrent
\[
  \SI{2259}{\kilo\joule} + \SI{2.6e2}{\kilo\joule} \approx \SI{2.5e3}{\kilo\joule}
\]
Per gram gir altså \doubleunderline{dampen nesten ti ganger så mye energi} til huden som vannet.

\smallskip\noindent
\textit{Merk:} Svaret gjelder for like stor \emph{masse}. Vanndamp har mye mindre massetetthet enn vann. En liter vanndamp har bare en masse på omtrent \SI{0.6}{\gram}, mens en liter vann har massen \SI{1}{\kilo\gram}. En liter kokende vann over hånda er derfor langt verre enn en liter damp. Men en dampstråle fra en kjele eller en dampkoker fortsetter å kondensere på huden så lenge den treffer, og da blir det fort mye energi.
}

\oppgaver{Cola og is}{k8oppg7}{%
%
En boks cola ($m_\text{c} = \SI{350}{\gram}$) har romtemperatur, $t_\text{c} = 20\,^\circ$C. Du legger i $m_\text{is} = \SI{20}{\gram}$ is med temperatur $t_\text{is} = -5.0\,^\circ$C. Anta at cola oppfører seg termisk som vann, og at det ikke utveksles varme med omgivelsene.
\begin{enumerate}[a)]
    \item Beregn energien som trengs for å varme opp isen til $0\,^\circ$C og smelte den.
    \item Finn slutt-temperaturen $t_\text{f}$ i blandingen.
\end{enumerate}
\textit{(Spesifikk varmekapasitet: $c_\text{vann} = c_\text{cola} = \SI{4.2}{kJ/(kg\,K)}$, $c_\text{is} = \SI{2.1}{kJ/(kg\,K)}$. Spesifikk smeltevarme: $\ell_\text{smelte} = \SI{334}{kJ/kg}$.)}
%
%
}

\svar[title={Løsningsforslag}]{%
\noindent
\textbf{a)} Først varmer vi isen fra $-5.0\,^\circ$C til $0\,^\circ$C, altså $\Delta T = \SI{5.0}{\kelvin}$:
\[
  Q_1 = c_\text{is}\, m_\text{is}\,\Delta T
      = \SI{2.1}{kJ/(kg\,K)}\cdot\SI{0.020}{\kilo\gram}\cdot\SI{5.0}{\kelvin}
      = \SI{0.21}{\kilo\joule}
\]
Så smelter vi isen ved $0\,^\circ$C:
\[
  Q_2 = \ell_\text{smelte}\, m_\text{is}
      = \SI{334}{kJ/kg}\cdot\SI{0.020}{\kilo\gram}
      = \SI{6.68}{\kilo\joule}
\]
Totalt:
\[
  Q_1 + Q_2 = \SI{0.21}{\kilo\joule} + \SI{6.68}{\kilo\joule} = \SI{6.89}{\kilo\joule} \approx \doubleunderline{\SI{6.9}{\kilo\joule}}
\]
Vi sjekker at all isen smelter. Colaen kan avgi høyst $4.2\cdot 0.350\cdot 20\,\si{\kilo\joule} = \SI{29.4}{\kilo\joule}$ før den når $0\,^\circ$C. Det er mye mer enn \SI{6.89}{\kilo\joule}, så all isen smelter, og $t_\text{f} > 0\,^\circ$C.

\smallskip\noindent
\textbf{b)} Varmen colaen avgir, er lik varmen isen (og smeltevannet) tar opp. Smeltevannet varmes fra $0\,^\circ$C til $t_\text{f}$:
\[
  c_\text{cola}\, m_\text{c}\,(20 - t_\text{f}) = Q_1 + Q_2 + c_\text{vann}\, m_\text{is}\,(t_\text{f} - 0)
\]
Med tallene inn (energi i kJ, temperatur i $^\circ$C):
\[
  4.2\cdot 0.350\cdot(20 - t_\text{f}) = 6.89 + 4.2\cdot 0.020\cdot t_\text{f}
\]
\[
  29.4 - 1.47\,t_\text{f} = 6.89 + 0.084\,t_\text{f}
  \quad\implies\quad
  t_\text{f} = \frac{29.4 - 6.89}{1.47 + 0.084} = \SI{14.5}{\celsius} \approx \doubleunderline{\SI{14}{\celsius}}
\]
}

\oppgaver{Vannkoker}{k8oppg8}{%
%
En vannkoker har effekten $P = \SI{2.0}{\kilo\watt}$. Du fyller den med \SI{1.5}{\litre} vann med temperatur $12\,^\circ$C. Anta at all energien fra vannkokeren går til vannet.
\begin{enumerate}[a)]
    \item Hvor mye energi trengs for å varme opp vannet til $100\,^\circ$C?
    \item Hvor lang tid tar det før vannet koker?
    \item Vannkokeren slår seg ikke av, og koker bort \SI{0.20}{\kilo\gram} vann før du slår den av. Hvor lenge sto den på etter at vannet begynte å koke?
\end{enumerate}
%
%
}

\svar[title={Løsningsforslag}]{%
\noindent
\textbf{a)} \SI{1.5}{\litre} vann har massen $m = \SI{1.5}{\kilo\gram}$. Temperaturendringen er $\Delta T = 100 - 12 = \SI{88}{\kelvin}$:
\[
  Q = c\, m\,\Delta T = \SI{4.2}{kJ/(kg\,K)}\cdot\SI{1.5}{\kilo\gram}\cdot\SI{88}{\kelvin}
    = \SI{554}{\kilo\joule} \approx \doubleunderline{\SI{5.5e2}{\kilo\joule}}
\]

\smallskip\noindent
\textbf{b)} Effekt er energi per tid, $P = E/t$:
\[
  t = \frac{Q}{P} = \frac{\SI{554e3}{\joule}}{\SI{2.0e3}{\watt}} = \SI{277}{\second} \approx \doubleunderline{\SI{2.8e2}{\second}}
\]
Det er omtrent \SI{4.6}{\minute}.

\smallskip\noindent
\textbf{c)} Energien som trengs for å fordampe \SI{0.20}{\kilo\gram} vann, er
\[
  Q = \ell_\text{fordamp}\, m = \SI{2259}{kJ/kg}\cdot\SI{0.20}{\kilo\gram} = \SI{452}{\kilo\joule}
\]
\[
  t = \frac{Q}{P} = \frac{\SI{452e3}{\joule}}{\SI{2.0e3}{\watt}} = \SI{226}{\second} \approx \doubleunderline{\SI{2.3e2}{\second}}
\]
Det er nesten fire minutter, selv om bare en liten del av vannet koker bort.
}

\oppgaver{Varme eller arbeid?}{k8oppg1}{%
%
Avgjør om energien overføres til systemet som \emph{varme} $Q$ eller som \emph{arbeid} $W$. Systemet står i parentes.
\begin{enumerate}[a)]
    \item Du gnir hendene hardt mot hverandre (hendene).
    \item En gryte med vann står på en varm kokeplate (vannet).
    \item Lufta i en sykkelpumpe blir presset sammen av stempelet (lufta).
    \item Sola varmer opp en stein på stranda (steinen).
    \item En elektrisk mikser rører i en bolle med kakedeig (deigen).
\end{enumerate}
%
%
}

\svar[title={Løsningsforslag}]{%
\noindent
Varme er energi som overføres \emph{bare} på grunn av en temperaturforskjell. All annen energioverføring, der noe skyver, drar eller rører, er arbeid.

\smallskip\noindent
\textbf{a)} \doubleunderline{Arbeid}. Friksjonskraften virker langs bevegelsen og gjør arbeid på hendene.

\smallskip\noindent
\textbf{b)} \doubleunderline{Varme}. Kokeplata er varmere enn gryta og vannet, og energien går fra det varme til det kalde.

\smallskip\noindent
\textbf{c)} \doubleunderline{Arbeid}. Stempelet skyver på lufta og presser den sammen.

\smallskip\noindent
\textbf{d)} \doubleunderline{Varme}. Sola er mye varmere enn steinen. Energien kommer som stråling, og det regner vi som varme.

\smallskip\noindent
\textbf{e)} \doubleunderline{Arbeid}. Mikseren skyver deigen rundt. Deigen blir litt varmere, selv om ingenting varmt er i kontakt med den.
}

\oppgaver{Sykkelpumpe}{k8oppg3}{%
%
En sykkelpumpe komprimerer luft raskt, slik at prosessen er tilnærmet adiabatisk. Et arbeid på $W = \SI{30}{\joule}$ utføres på lufta.
\begin{enumerate}[a)]
    \item Hva er endringen i den indre energien til lufta?
    \item Stiger eller synker temperaturen i pumpa? Forklar hvorfor.
\end{enumerate}
%
%
}

\svar[title={Løsningsforslag}]{%
\noindent
\textbf{a)} Prosessen er adiabatisk, så $Q = 0$. Termodynamikkens 1.\ lov gir
\[
  \Delta U = Q + W = 0 + \SI{30}{\joule} = \doubleunderline{+\SI{30}{\joule}}
\]

\smallskip\noindent
\textbf{b)} $\Delta U > 0$, så den indre energien øker. Den indre energien til lufta er knyttet til bevegelsen til molekylene, og dermed til temperaturen. Temperaturen \doubleunderline{stiger}. Stempelet beveger seg mot molekylene og gir dem større fart når de spretter mot det. Du kan kjenne at pumpa blir varm etter bruk.
}

\oppgaver{Gass som komprimeres}{k8oppg4}{%
%
Et stempel presser sammen en gass i en sylinder og gjør et arbeid på \SI{250}{\joule} på gassen. Samtidig avgir gassen \SI{80}{\joule} varme til omgivelsene. Anta at den indre energien til gassen bare avhenger av temperaturen.
\begin{enumerate}[a)]
    \item Finn endringen $\Delta U$ i den indre energien til gassen.
    \item Stiger eller synker temperaturen til gassen?
    \item Hvor mye varme måtte gassen ha avgitt for at temperaturen skulle vært uendret?
    \item Senere tilføres gassen \SI{600}{\joule} varme. Gassen utvider seg og løfter stempelet, som har masse \SI{25}{\kilo\gram}, rett opp. Den indre energien øker med \SI{450}{\joule}. Hvor høyt løftes stempelet? Se bort fra lufttrykket over stempelet og friksjon.
\end{enumerate}
%
%
}

\svar[title={Løsningsforslag}]{%
\noindent
\textbf{a)} Arbeidet utføres \emph{på} gassen, så $W = +\SI{250}{\joule}$. Gassen \emph{avgir} varme, så $Q = -\SI{80}{\joule}$. Termodynamikkens 1.\ lov gir
\[
  \Delta U = Q + W = -\SI{80}{\joule} + \SI{250}{\joule} = \doubleunderline{+\SI{170}{\joule}}
\]

\smallskip\noindent
\textbf{b)} Den indre energien øker. Siden den indre energien bare avhenger av temperaturen, \doubleunderline{stiger temperaturen}.

\smallskip\noindent
\textbf{c)} Temperaturen er uendret når $\Delta U = 0$:
\[
  0 = Q + W \quad\implies\quad Q = -W = \doubleunderline{-\SI{250}{\joule}}
\]
Gassen måtte altså ha avgitt \SI{250}{\joule} varme, like mye som arbeidet stempelet gjør på den.

\smallskip\noindent
\textbf{d)} Nå er $Q = +\SI{600}{\joule}$ og $\Delta U = +\SI{450}{\joule}$. Termodynamikkens 1.\ lov gir
\[
  W = \Delta U - Q = \SI{450}{\joule} - \SI{600}{\joule} = -\SI{150}{\joule}
\]
Det negative fortegnet betyr at gassen gjør \SI{150}{\joule} arbeid på stempelet. Arbeidet går til å løfte stempelet mot tyngdekraften, $W_\text{gass} = m\mathrm{g}h$:
\[
  h = \frac{W_\text{gass}}{m\mathrm{g}} = \frac{\SI{150}{\joule}}{\SI{25}{\kilo\gram}\cdot\SI{9.81}{\metre\per\second\squared}} = \SI{0.612}{\metre} \approx \doubleunderline{\SI{0.61}{\metre}}
\]
}

\oppgaver{Sylindereksperiment}{k8oppg13}{%
%
En gass i en sylinder har volum $V_1 = \SI{30}{\deci\metre\cubed}$, trykk $p_1 = \SI{100}{\kilo\pascal}$ og temperatur $T_1 = \SI{300}{\kelvin}$. Stempelet beveger seg slik at gassen utvider seg ved konstant trykk (isobart), helt til temperaturen er $T_2 = \SI{600}{\kelvin}$.
\begin{enumerate}[a)]
    \item Hva er det nye volumet $V_2$?
    \item Hva skjer med den indre energien til gassen? Må gassen få tilført varme?
\end{enumerate}
%
%
}

\svar[title={Løsningsforslag}]{%
\noindent
\textbf{a)} Antall molekyler og trykket er konstant. Tilstandslikningen $pV = Nk_\text{B}T$ gir da at $V/T$ er konstant:
\[
  \frac{V_1}{T_1} = \frac{V_2}{T_2}
  \quad\implies\quad
  V_2 = V_1\cdot\frac{T_2}{T_1} = \SI{30}{\deci\metre\cubed}\cdot\frac{\SI{600}{\kelvin}}{\SI{300}{\kelvin}} = \doubleunderline{\SI{60}{\deci\metre\cubed}}
\]

\smallskip\noindent
\textbf{b)} Den absolutte temperaturen dobles, og dermed dobles også den gjennomsnittlige kinetiske energien til molekylene. Den indre energien \doubleunderline{øker}, $\Delta U > 0$.

Samtidig utvider gassen seg og skyver stempelet ut. Gassen gjør altså arbeid på omgivelsene, og $W < 0$ i termodynamikkens 1.\ lov. Da gir
\[
  Q = \Delta U - W
\]
at $Q$ er positiv og \emph{større} enn $\Delta U$. Gassen \doubleunderline{må få tilført varme}: noe av varmen øker den indre energien, og resten går med til arbeidet på stempelet.
}

\oppgaver{Bildekk om vinteren}{k8oppg14}{%
%
Et bildekk blir fylt med luft i en garasje som holder $20\,^\circ$C. Det absolutte trykket i dekket er da \SI{3.0e5}{\pascal}, og volumet av lufta i dekket er \SI{0.030}{\metre\cubed}. Så kjøres bilen ut i kulda, der det er $-15\,^\circ$C. Anta at volumet av dekket ikke endres.
\begin{enumerate}[a)]
    \item Hvor mange luftmolekyler er det i dekket?
    \item Hva blir det absolutte trykket i dekket når lufta har fått samme temperatur som omgivelsene?
    \item En dekktrykkmåler viser \emph{overtrykket}, altså forskjellen mellom trykket i dekket og lufttrykket utenfor, $p_0 = \SI{1.0e5}{\pascal}$. Hva viser måleren før og etter? Hvor mange prosent har overtrykket sunket?
\end{enumerate}
%
%
}

\svar[title={Løsningsforslag}]{%
\noindent
Temperaturene i kelvin er $T_1 = 20 + 273 = \SI{293}{\kelvin}$ og $T_2 = -15 + 273 = \SI{258}{\kelvin}$.

\smallskip\noindent
\textbf{a)} Vi løser tilstandslikningen for $N$:
\[
  N = \frac{p_1V}{k_\text{B}T_1} = \frac{\SI{3.0e5}{\pascal}\cdot\SI{0.030}{\metre\cubed}}{\SI{1.38e-23}{\joule\per\kelvin}\cdot\SI{293}{\kelvin}} = \doubleunderline{\num{2.2e24}}
\]

\smallskip\noindent
\textbf{b)} $N$ og $V$ er konstante, så $p/T$ er konstant:
\[
  \frac{p_1}{T_1} = \frac{p_2}{T_2}
  \quad\implies\quad
  p_2 = p_1\cdot\frac{T_2}{T_1} = \SI{3.0e5}{\pascal}\cdot\frac{\SI{258}{\kelvin}}{\SI{293}{\kelvin}} = \SI{2.64e5}{\pascal} \approx \doubleunderline{\SI{2.6e5}{\pascal}}
\]

\smallskip\noindent
\textbf{c)} Før: $p_1 - p_0 = \SI{3.0e5}{\pascal} - \SI{1.0e5}{\pascal} = \doubleunderline{\SI{2.0e5}{\pascal}}$.

Etter: $p_2 - p_0 = \SI{2.64e5}{\pascal} - \SI{1.0e5}{\pascal} = \SI{1.64e5}{\pascal} \approx \doubleunderline{\SI{1.6e5}{\pascal}}$.

Overtrykket har sunket med
\[
  \frac{\SI{2.0e5}{\pascal} - \SI{1.64e5}{\pascal}}{\SI{2.0e5}{\pascal}} = 0.18 = \doubleunderline{\SI{18}{\percent}}
\]
Det absolutte trykket sank bare med omtrent \SI{12}{\percent}, men måleren viser overtrykket, og det synker mer. Derfor bør man sjekke dekktrykket når det blir kaldt.
}

\oppgaver{Arbeid i et pV-diagram}{k8oppg15}{%
%
En gass utvider seg ved konstant trykk $p = \SI{150}{\kilo\pascal}$, fra volumet \SI{2.0}{\litre} til \SI{5.0}{\litre}.
\begin{enumerate}[a)]
    \item Tegn prosessen i et $pV$-diagram.
    \item Finn arbeidet gassen gjør på omgivelsene. Hva blir $W$ i termodynamikkens 1.\ lov?
    \item Under prosessen blir det tilført \SI{1.1}{\kilo\joule} varme til gassen. Finn endringen i den indre energien til gassen.
\end{enumerate}
%
%
}

\svar[title={Løsningsforslag}]{%
\noindent
\textbf{a)} Trykket er konstant, så prosessen er en horisontal linje fra \SI{2.0}{\litre} til \SI{5.0}{\litre}. Arealet under linjen er arbeidet gassen gjør.

\begin{center}
\begin{tikzpicture}[scale=0.9, font=\small, >={Stealth[length=6pt,width=4pt]}]
  \fill[defscol!15] (2.0,0) rectangle (5.0,2.5);
  \draw[->] (0,0) -- (6.2,0) node[right] {$V$/L};
  \draw[->] (0,0) -- (0,3.4) node[above] {$p$/kPa};
  \draw[ultra thick, defscol, ->] (2.0,2.5) -- (5.0,2.5);
  \fill (2.0,2.5) circle (2pt) (5.0,2.5) circle (2pt);
  \draw[dashed, gray] (0,2.5) -- (2.0,2.5);
  \draw[dashed, gray] (2.0,0) -- (2.0,2.5) (5.0,0) -- (5.0,2.5);
  \node[left] at (0,2.5) {150};
  \node[below] at (2.0,0) {2.0};
  \node[below] at (5.0,0) {5.0};
  \node[defscol!60!black] at (3.5,1.25) {areal $= p\,\Delta V$};
\end{tikzpicture}
\end{center}

\noindent
\textbf{b)} Arbeidet gassen gjør, er arealet under linjen. Med $\Delta V = \SI{5.0}{\litre} - \SI{2.0}{\litre} = \SI{3.0}{\litre} = \SI{3.0e-3}{\metre\cubed}$:
\[
  W_\text{gass} = p\,\Delta V = \SI{150e3}{\pascal}\cdot\SI{3.0e-3}{\metre\cubed} = \doubleunderline{\SI{4.5e2}{\joule}}
\]
I 1.\ lov er $W$ arbeidet som gjøres \emph{på} gassen. Gassen gjør arbeid på omgivelsene, så $W = \doubleunderline{-\SI{4.5e2}{\joule}}$.

\smallskip\noindent
\textbf{c)}
\[
  \Delta U = Q + W = \SI{1.1e3}{\joule} - \SI{4.5e2}{\joule} = \doubleunderline{\SI{6.5e2}{\joule}}
\]
}

\oppgaver{To prosesser etter hverandre}{k8oppg16}{%
%
En gass er i tilstand A med trykk $p_\text{A} = \SI{100}{\kilo\pascal}$, volum $V_\text{A} = \SI{4.0}{\litre}$ og temperatur $T_\text{A} = \SI{300}{\kelvin}$. Gassen blir først presset sammen isotermt til volumet \SI{1.0}{\litre} (tilstand B). Deretter blir den varmet opp ved konstant volum til \SI{600}{\kelvin} (tilstand C).
\begin{enumerate}[a)]
    \item Finn trykket i tilstand B.
    \item Finn trykket i tilstand C.
    \item Skisser de to prosessene i et $pV$-diagram.
    \item Gjør gassen noe arbeid fra B til C? Begrunn svaret.
\end{enumerate}
%
%
}

\svar[title={Løsningsforslag}]{%
\noindent
\textbf{a)} A til B er isoterm ($T$ konstant), så $pV$ er konstant:
\[
  p_\text{A}V_\text{A} = p_\text{B}V_\text{B}
  \quad\implies\quad
  p_\text{B} = p_\text{A}\cdot\frac{V_\text{A}}{V_\text{B}} = \SI{100}{\kilo\pascal}\cdot\frac{\SI{4.0}{\litre}}{\SI{1.0}{\litre}} = \doubleunderline{\SI{4.0e2}{\kilo\pascal}}
\]

\smallskip\noindent
\textbf{b)} B til C er isokor ($V$ konstant), så $p/T$ er konstant. Temperaturen i B er fortsatt \SI{300}{\kelvin}:
\[
  p_\text{C} = p_\text{B}\cdot\frac{T_\text{C}}{T_\text{B}} = \SI{4.0e2}{\kilo\pascal}\cdot\frac{\SI{600}{\kelvin}}{\SI{300}{\kelvin}} = \doubleunderline{\SI{8.0e2}{\kilo\pascal}}
\]

\smallskip\noindent
\textbf{c)} Fra A til B følger gassen en isoterm, en kurve der $pV$ er konstant. Fra B til C går gassen rett opp, fordi volumet er konstant.

\begin{center}
\begin{tikzpicture}[xscale=1.1, yscale=0.42, font=\small, >={Stealth[length=6pt,width=4pt]}]
  \draw[->] (0,0) -- (5,0) node[right] {$V$/L};
  \draw[->] (0,0) -- (0,9.2) node[above] {$p$/kPa};
  \draw[ultra thick, headCol, domain=4:1, samples=40, ->] plot (\x, {4/\x});
  \draw[ultra thick, red!70!black, ->] (1,4) -- (1,8);
  \fill (4,1) circle (3pt) node[above right] {A};
  \fill (1,4) circle (3pt) node[left] {B};
  \fill (1,8) circle (3pt) node[left] {C};
  \foreach \x in {1,2,3,4} \draw (\x,0.15) -- (\x,-0.15) node[below] {\x};
  \foreach \y/\l in {1/100,4/400,8/800} \draw (0.06,\y) -- (-0.06,\y) node[left] {\l};
  \node[headCol!80!black, font=\footnotesize] at (2.9,2.3) {isoterm};
  \node[red!70!black, font=\footnotesize, right] at (1.05,6) {isokor};
\end{tikzpicture}
\end{center}

\noindent
\textbf{d)} \doubleunderline{Nei}. Fra B til C er volumet konstant. Stempelet beveger seg ikke, og arealet under linjen i $pV$-diagrammet er null. Da er $W = 0$, og all varmen som tilføres, går til å øke den indre energien.
}

\oppgaver{Bremser som blir varme}{k8oppg17}{%
%
En bil med masse \SI{1.2e3}{\kilo\gram} kjører i \SI{72}{\kilo\metre\per\hour} og bremser helt ned til stillstand. Anta at all den kinetiske energien blir til indre energi i de fire bremseskivene av jern, som har en samlet masse på \SI{12}{\kilo\gram}.
\begin{enumerate}[a)]
    \item Finn den kinetiske energien til bilen før bremsingen.
    \item Hvor mye stiger temperaturen i bremseskivene? \textit{($c_\text{jern} = \SI{0.45}{kJ/(kg\,K)}$.)}
\end{enumerate}
%
%
}

\svar[title={Løsningsforslag}]{%
\noindent
\textbf{a)} Farten er $v = \SI{72}{\kilo\metre\per\hour} = \frac{72}{3.6}\,\si{\metre\per\second} = \SI{20}{\metre\per\second}$:
\[
  E_\text{K} = \frac{1}{2}mv^2 = \frac{1}{2}\cdot\SI{1.2e3}{\kilo\gram}\cdot(\SI{20}{\metre\per\second})^2 = \doubleunderline{\SI{2.4e5}{\joule}}
\]

\smallskip\noindent
\textbf{b)} All den kinetiske energien blir til indre energi i bremseskivene, $Q = c\,m\,\Delta T = E_\text{K}$:
\[
  \Delta T = \frac{E_\text{K}}{c\,m} = \frac{\SI{2.4e5}{\joule}}{\SI{450}{\joule\per(\kilo\gram\kelvin)}\cdot\SI{12}{\kilo\gram}} = \SI{44.4}{\kelvin} \approx \doubleunderline{\SI{44}{\kelvin}}
\]
Strengt tatt er dette \emph{arbeid} (friksjon), ikke varme, men regningen blir den samme.
}

\oppgaver{Isbiter i te}{k8oppg18}{%
%
Du har \SI{0.30}{\kilo\gram} te med temperatur $80\,^\circ$C, men vil drikke den ved $50\,^\circ$C. Du legger i is med temperatur $0\,^\circ$C. Hvor mye is må du bruke? Anta at teen oppfører seg termisk som vann, og se bort fra koppen og varmetap til omgivelsene.
\textit{($c_\text{vann} = \SI{4.2}{kJ/(kg\,K)}$, $\ell_\text{smelte} = \SI{334}{kJ/kg}$.)}
%
%
}

\svar[title={Løsningsforslag}]{%
\noindent
Teen kjøles fra $80\,^\circ$C til $50\,^\circ$C og avgir varmen
\[
  Q_\text{te} = c_\text{vann}\,m_\text{te}\,\Delta T = \SI{4.2}{kJ/(kg\,K)}\cdot\SI{0.30}{\kilo\gram}\cdot\SI{30}{\kelvin} = \SI{37.8}{\kilo\joule}
\]
Isen med masse $m$ må først smelte ved $0\,^\circ$C, og så må smeltevannet varmes opp til $50\,^\circ$C:
\[
  Q_\text{is} = \ell_\text{smelte}\,m + c_\text{vann}\,m\cdot\SI{50}{\kelvin}
  = m\,(\SI{334}{kJ/kg} + \SI{210}{kJ/kg}) = m\cdot\SI{544}{kJ/kg}
\]
Varmen teen avgir, er lik varmen isen tar opp:
\[
  m = \frac{\SI{37.8}{\kilo\joule}}{\SI{544}{kJ/kg}} = \SI{0.0695}{\kilo\gram} \approx \doubleunderline{\SI{69}{\gram}}
\]
}

\oppgaver{Ballong i fryseren}{k8oppg19}{%
%
En ballong har volumet \SI{3.0}{\litre} ved $22\,^\circ$C. Den legges i en fryser som holder $-18\,^\circ$C. Anta at trykket i ballongen er konstant. Hva blir volumet til ballongen når lufta inni har fått samme temperatur som fryseren?
%
%
}

\svar[title={Løsningsforslag}]{%
\noindent
Temperaturene i kelvin er $T_1 = 22 + 273 = \SI{295}{\kelvin}$ og $T_2 = -18 + 273 = \SI{255}{\kelvin}$. Trykket og antall molekyler er konstant, så $V/T$ er konstant:
\[
  V_2 = V_1\cdot\frac{T_2}{T_1} = \SI{3.0}{\litre}\cdot\frac{\SI{255}{\kelvin}}{\SI{295}{\kelvin}} = \SI{2.59}{\litre} \approx \doubleunderline{\SI{2.6}{\litre}}
\]
Ballongen krymper litt, men ikke så mye som man kanskje skulle tro. Det er temperaturen i kelvin som teller, og den synker bare med omtrent \SI{14}{\percent}.
}

\oppgaver{Molekyler i et klasserom}{k8oppg20}{%
%
Et klasserom er \SI{8.0}{\metre} langt, \SI{6.0}{\metre} bredt og \SI{3.0}{\metre} høyt. Lufttrykket er \SI{1.0e5}{\pascal}, og temperaturen er $20\,^\circ$C.
\begin{enumerate}[a)]
    \item Hvor mange luftmolekyler er det i rommet?
    \item Temperaturen stiger til $25\,^\circ$C, mens trykket er det samme fordi døra står på gløtt. Omtrent hvor mange molekyler har forlatt rommet?
\end{enumerate}
%
%
}

\svar[title={Løsningsforslag}]{%
\noindent
\textbf{a)} Volumet er $V = \SI{8.0}{\metre}\cdot\SI{6.0}{\metre}\cdot\SI{3.0}{\metre} = \SI{144}{\metre\cubed}$, og $T_1 = 20 + 273 = \SI{293}{\kelvin}$:
\[
  N_1 = \frac{pV}{k_\text{B}T_1} = \frac{\SI{1.0e5}{\pascal}\cdot\SI{144}{\metre\cubed}}{\SI{1.38e-23}{\joule\per\kelvin}\cdot\SI{293}{\kelvin}} = \SI{3.56e27}{} \approx \doubleunderline{\num{3.6e27}}
\]

\smallskip\noindent
\textbf{b)} Ved $T_2 = 25 + 273 = \SI{298}{\kelvin}$, med samme $p$ og $V$:
\[
  N_2 = \frac{pV}{k_\text{B}T_2} = N_1\cdot\frac{T_1}{T_2} = \num{3.56e27}\cdot\frac{293}{298} = \num{3.50e27}
\]
Antall molekyler som har forlatt rommet:
\[
  N_1 - N_2 = \num{3.56e27} - \num{3.50e27} \approx \doubleunderline{\num{6.0e25}}
\]
Det er bare omtrent \SI{1.7}{\percent} av molekylene, men likevel et enormt tall.
}

\oppgaver{Kobber i vann}{k8oppg21}{%
%
Et lodd av kobber med masse \SI{0.20}{\kilo\gram} blir varmet opp til $100\,^\circ$C i kokende vann. Så slippes det ned i en godt isolert beholder med \SI{0.50}{\kilo\gram} vann med temperatur $15\,^\circ$C. Se bort fra beholderen. Hva blir slutt-temperaturen?
\textit{($c_\text{kobber} = \SI{0.39}{kJ/(kg\,K)}$, $c_\text{vann} = \SI{4.2}{kJ/(kg\,K)}$.)}
%
%
}

\svar[title={Løsningsforslag}]{%
\noindent
Kobberet avgir varme og kjøles fra $100\,^\circ$C til slutt-temperaturen $t$. Vannet tar opp varmen og varmes fra $15\,^\circ$C til $t$. Varmen kobberet avgir, er lik varmen vannet tar opp:
\[
  c_\text{k}\,m_\text{k}\,(100 - t) = c_\text{v}\,m_\text{v}\,(t - 15)
\]
Med tallene inn (i kJ/K):
\[
  0.39\cdot 0.20\cdot(100 - t) = 4.2\cdot 0.50\cdot(t - 15)
\]
\[
  7.8 - 0.078\,t = 2.1\,t - 31.5
  \quad\implies\quad
  t = \frac{7.8 + 31.5}{2.1 + 0.078} = \SI{18.0}{\celsius} \approx \doubleunderline{\SI{18}{\celsius}}
\]
Vannet blir bare litt varmere. Vann har mye større spesifikk varmekapasitet enn kobber, og det er i tillegg mer vann enn kobber.
}
